% ============================================================
% Section 168: 碱金属价电子受屏蔽势作用时的能级
% ============================================================
\section{量子力学：碱金属价电子受屏蔽势作用时的能级}
\label{sec:alkali-screened-potential}

\begin{modelbox}[题目：碱金属价电子受屏蔽势作用时的能级]
设碱金属原子中的价电子所受原子实（原子核 $+$ 满壳电子）的作用近似表示为
\[
V(r)=-\frac{e^2}{r}-\lambda\frac{e^2a}{r^2},\quad 0<\lambda\ll1,
\]
其中 $a=\hbar^2/(\mu e^2)$ 为玻尔半径，$\mu$ 为约化质量。试求价电子的能级，并与氢原子能级作比较。
\end{modelbox}

\begin{tipbox}[考点快速分析]
\begin{itemize}
    \item \textbf{有效角动量}：将 $1/r^2$ 项并入离心势，引入有效角量子数 $l'$
    \item \textbf{类氢映射}：方程化为标准类氢形式，直接借用氢原子能级公式
    \item \textbf{量子亏损}：$\Delta_l=l'-l\approx-\lambda/(l+\frac12)$，消除 $l$ 简并
    \item \textbf{物理图像}：低 $l$ 态贯穿原子实更深，屏蔽效应更强，能级下移更多
\end{itemize}
\end{tipbox}

\begin{derivationbox}[标准解题过程]
\textbf{1. 径向方程与有效角动量}

有心力场径向方程（$u(r)=rR(r)$）：
\[
-\frac{\hbar^2}{2\mu}\frac{d^2u}{dr^2}+\Bigl[-\frac{e^2}{r}+\frac{l(l+1)\hbar^2}{2\mu r^2}-\lambda\frac{e^2a}{r^2}\Bigr]u=Eu.
\]

利用 $a=\hbar^2/(\mu e^2)$，附加项系数化为 $2\mu\lambda e^2a/\hbar^2=2\lambda$。于是离心势变为
\[
\frac{l(l+1)-2\lambda}{r^2}.
\]

引入有效角量子数 $l'$ 使 $l'(l'+1)=l(l+1)-2\lambda$，解得
\begin{equation}
\boxed{l'=\sqrt{\Bigl(l+\frac12\Bigr)^2-2\lambda}-\frac12\approx l-\frac{\lambda}{l+\frac12}\quad(\lambda\ll1).}
\end{equation}

\textbf{2. 能级与波函数}

方程化为标准类氢形式，仅需将 $l\to l'$。束缚态能级为
\begin{equation}
\boxed{E_{n_rl}=-\frac{\mu e^4}{2\hbar^2}\frac{1}{(n')^2}=-\frac{e^2}{2a}\frac{1}{(n_r+l'+1)^2},}
\end{equation}
其中 $n_r=0,1,2,\ldots$ 为径向量子数，$n'=n_r+l'+1$ 为有效主量子数（一般非整数）。

径向波函数为
\[
R_{n'l'}(r)=N_{n'l'}\,e^{-r/(n'a)}\Bigl(\frac{2r}{n'a}\Bigr)^{l'}F\Bigl(-n_r,\,2l'+2,\,\frac{2r}{n'a}\Bigr).
\]

\textbf{3. 量子亏损与物理讨论}

定义量子亏损 $\Delta_l=l'-l\approx-\lambda/(l+\frac12)\ll0$。能级可写为
\[
E_{nl}\approx-\frac{\mu e^4}{2\hbar^2}\frac{1}{\bigl(n-\frac{\lambda}{l+1/2}\bigr)^2}.
\]

与氢原子比较：
\begin{itemize}
    \item 氢原子能级仅依赖 $n$，高度 $l$ 简并；
    \item 碱金属中屏蔽势破坏纯粹库仑场，能级依赖 $l$，消除了 $l$ 简并；
    \item $l$ 越小，$|\Delta_l|$ 越大，能级越低（电子轨道越扁，贯穿原子实越深，有效核电荷越大）；
    \item 高 $l$ 态 $|\Delta_l|\to0$，趋近氢原子能级。
\end{itemize}
这与碱金属光谱（如钠 D 线 $3p\to3s$）的实验规律完全一致。
\end{derivationbox}

\begin{formulabox}[答案汇总]
\begin{center}
\begin{tabular}{ll}
\toprule
\textbf{结论} & \textbf{结果} \\
\midrule
有效角量子数 & $l'=\sqrt{(l+\frac12)^2-2\lambda}-\frac12$ \\
能级公式 & $E_{n_rl}=-\dfrac{\mu e^4}{2\hbar^2}\dfrac{1}{(n_r+l'+1)^2}$ \\
量子亏损（$\lambda\ll1$） & $\Delta_l\approx-\dfrac{\lambda}{l+\frac12}$ \\
\bottomrule
\end{tabular}
\end{center}
\end{formulabox}

\begin{insightbox}[物理图像]
\begin{itemize}
    \item 碱金属光谱的量子亏损是原子物理中的经典范例。实验上通过测量谱线可反推出 $\Delta_l$，进而研究原子实的极化和贯穿效应。
    \item 本题的方法——将微扰项吸收进有效量子数——在量子力学中十分通用，类似技巧也出现在离心势修正、自旋-轨道耦合等问题中。
    \item 对于多电子原子，屏蔽效应更复杂，但量子亏损概念仍适用，可推广为量子亏损理论（QDT）。
\end{itemize}
\end{insightbox}
